\chapter{The Scissors Crisis}

The ``link with the peasant'' which NEP was designed to establish remained the watchword of Soviety policy for several years. Few doubted its necessity. ``Only an agreement with the peasantry'', Lenin had said at the tenth congress, ``can save the socialist revolution in Russia until the revolution has occurred in other countries''. When, after the terrible famine of 1921-1922, agriculture quickly revived, and recovery began to spread to other sectors of the economy, NEP had been triumphantly vindicated. Yet, once the danger was over, and memories of the privations of war communism receded into the past, the mood of relief and acquiescence slowly faded, and was overtaken by a sense of uneasiness at so radical a departure from the hopes and expectations of an advance into socialism which had inspired the early triumphs of the revolution. In the long run someone carried the cost of concessions made to the peasant; and some consequences of NEP, direct or indirect, were unlooked-for and unwelcome. In little more than two years the country was in the throes of a fresh crisis which, though less dramatic than the crisis preceding the introduction of NEP, deeply affected every sector of a now expanding economy. 
% 注：去除了 "llbegan"、"^'" 等边缘乱码，修正了 "sociahsrn" 的 OCR 拼写错误为 "socialism"。

The impact of NEP on industry was less direct than on agriculture, and mainly negative. Its first effect was to stimulate the recovery of rural and artisan industries, both because they had suffered less than factory industry in the civil war, and could more easily be brought back into production, and because they were the main suppliers of the simple consumer goods which the peasant wanted to purchase with the proceeds of his sale of agricultural products. The campaign for the nationalization of industry was halted. Large-scale industry (Lenin's ``commanding heights'') remained in state hands, but with two important modifications. In the first place, a large measure of devolution was practised. State industry was divided into three categories, ``Union'', ``republican'' and ``local''. ``Union'' industry was administered by the Vesenkha of the USSR, ``republican'' industry by the Vesenkhas of the constituent republics; and, within the republics, provinces and regions set up their own Councils of National Economy (Sovnarkhozy) which were responsible for ``local'' industry. A fluctuating degree of supervision was exercised by the higher over the lower organs. But practical considerations dictated a large measure of autonomy. At lower levels private industry was encouraged. Enterprises employing less than 20 workers were exempt from nationalization. Larger enterprises which had already been taken over might be leased back to individual entrepreneurs, often to their original owners. Rural, artisan and cooperative industries operated and expanded with official approval. 
% 注：去除了 "I", "A", "1 f" 等断行乱码，修正 "Agriculture" 为小写 "agriculture" 以符合语法。

Secondly, the direct administration of factory industry by Vesenkha through its glavki and centres was abolished. Industries were organized in trusts, which operated a group of enterprises as a single entity; the average number of enterprises in a trust was ten. The largest trusts were in the textile and metallurgical industries; the largest textile trust employed over 50,000 workers. The essential feature of the trusts was that they were no longer financed from the state budget, but were instructed to work on principles of commercial accounting (khozraschet), and to earn profits which were paid, subject to some deductions, to the state as the owner of the fixed capital of the enterprises. Some essential industries were still obliged to deliver a proportion of their output to state institutions. Otherwise industry, like the peasant, was free to sell its products on the market at whatever price they would fetch. These arrangements were consonant with the spirit of NEP. But they were criticized in some party circles; a blunt instruction from Vesenkha to the trusts in 1923 to earn ``maximum profits'' excited unfavourable publicity. 
% 注：去除了 "lover" 之前的 "l" 识别错误。

A year after the introduction of NEP the stimulus given by it to the availability and circulation of commodities of all kinds could be regarded with some degree of complacency. Lenin had been aware of the dangers of ``freedom of trade'', which, as he said at the tenth congress, ``inevitably leads to the victory of capital, to its full restoration''. He seems at first to have envisaged the exchange of goods between town and country as a grandiose system of organized barter. But, as he later admitted, ``the exchange of goods broke loose'', and ``turned into buying and selling''; and he shocked some party stalwarts by telling his hearers to ``learn to trade''. In 1922 a Commercial Bourse was established in Moscow. The intention was doubtless to exercise some form of public control over the processes of trade. The result was to facilitate the operations of a new class of merchants who quickly came to be characterized as ``Nepmen''. Petty private trade had never been completely extinguished, even under war communism; the famous Sukharevka market in Moscow was a known, and tolerated, abuse. The rising class of Nepmen were no longer petty traders, but large-scale commercial entrepreneurs who spread their tentacles into every sector of the economy. The big industrial trusts were still able to control the wholesale market in their products. Retail shops known as GUM (State Universal Store) were opened under Vesenkha auspices in Moscow and some other cities. But they were not at first very successful; and existing consumer cooperatives made little progress. Retail trade everywhere was dominated and fostered by Nepmen. As trade began to flow in increasing abundance, an air of prosperity returned to the well-to-do quarters of the capital. Many once-familiar features, banished by the revolution, now re-appeared in the landscape. Krasin on a visit to the city in September 1922 wrote to his wife that ``Moscow looks all right, in some parts as it was before the war''. Foreign visitors commented, sourly or triumphantly, according to their bent, on the revival of such ``capitalist'' phenomena as prostitutes on the streets, and subservient waiters and cabmen soliciting tips. For the beneficiaries of NEP prospects looked rosy. The worst seemed over. The shortages and tensions of war communism had been relaxed. Recovery was on its way. 
% 注：去除了 "[by", "Ull", "^1", "once-famihar" (修正为 once-familiar)。

Before long, however, the deeper implications of NEP revealed themselves in several interconnected crises. The first was a crisis of prices. Now that the controls of war communism had been lifted, prices fluctuated wildly. A prices committee appointed in August 1921, and a commission for internal trade set up in May 1922, proved wholly ineffective. The hunger of the cities for agricultural products outstripped the hunger of the peasant for the products of industry, so that agricultural prices at first soared in relation to industrial prices. Industry, denuded of working capital and deprived of sources of credit, could finance itself only by selling its products on a falling market, thus depressing industrial prices still further. This process, which reached its height in the summer of 1922, resulted in a crisis of labour. Under war communism labour, like every commodity, was scarce, and unemployment unthinkable. Compulsory labour service had the advantage of ensuring food rations for mobilized workers. Now compulsory labour, except in the penal labour camps, was gradually abandoned, and free employment for wages returned; collective contracts began to be negotiated by trade unions on behalf of their members. But the number of jobs was now smaller than the number of workers seeking them. Employers continued for a long time to supply food rations to their workers; but these were now payments in kind calculated at market prices in lieu of wages. The vagaries of the price-index made wage-rates a constant subject of haggling, in which the worker had a poor bargaining position. Wages frequently fell into arrears owing to the inability of enterprises to find the cash to pay them. 
% 注：清除了 "jl Before", "jlrevealed", "Iwas", "(U. oilh", "ihunger", "{prices", "(was" 等边缘乱码。

The status of the trade unions was governed by the rather hollow compromise reached at the tenth party congress in March 1921, the limitations of which were shown up at the trade union congress two months later. Tomsky, having failed to resist an attempt to re-open the issues decided at the party congress, was severely reprimanded, relieved at the behest of the party authorities of his post as president of the trade union central council, and despatched on a mission to Central Asia. It may have been significant that Tomsky was succeeded as president by Andreev, originally a supporter of Trotsky's trade union platform. But this did not restore peace in the unions. In January 1922 the Politburo once more intervened with a resolution which recognized the existence of ``a series of contradictions between different tasks of the trade unions''---notably a contradiction between ``the defence of the interests of the toiling masses'' and the role of the unions as ``sharers of state power and builders of the national economy as a whole''. This formula seems to have paved the way for a reprieve of Tomsky, who at the next trade union congress in September 1922 was reinstated in his former post as president. The congress once more attempted to define the role of the unions. It was their function ``unconditionally to protect the interests of the workers''. On the other hand, the obligation rested on them to maintain and improve productivity, seen as the contribution of the workers to the building of a socialist order; and, though strikes were not formally prohibited, the right way to settle disputes was by negotiations between the trade union and the employer or economic administration concerned. It is noticeable that no distinction of substance appears to have been drawn at any time between the role of the trade unions in state and in private enterprises. Both contributed to essential production; and it was important that this process should not be disrupted. 
% 注：去除了 "/tasks", "/rthe", "^of", "TiTn", "ole" (修正为 role) 等乱码断字。

Discontent among the workers was fanned by the rising status and influence of the so-called ``Red managers''. In the civil war, former Tsarist officers had been employed to rebuild and command the Red Army; so also, in order to revive essential industries, former factory managers, and sometimes factory owners, had been pressed into service, in the guise of ``specialists'', as managers of nationalized enterprises, sometimes under the supervision of party members or workers. The system met the need for managerial skills, and was standardized and extended under NEP, when autonomous trusts and syndicates took the place of the glavki and centres of war communism. The group of ``Red managers'', in spite of their predominantly bourgeois origin and affiliations, acquired a recognized and respected place in the Soviet hierarchy; some of them were admitted to party membership---a reward for distinguished service. They received special rates of remuneration outside the normal wage-scales, and far in excess of them; and they exercised an increasingly powerful voice in industrial administration and industrial policy. Frequent and not unjustified accusations of their brutal and dictatorial attitude to the workers, reminiscent of the methods of the old regime, were symptomatic of the jealousies and resentments provoked by this apparent reversal of everything that the revolution had stood for in the factories. 
% 注：去除了 "I pf", "llexcess", "'/voice", "lof", "/resentments", "-^Athat" 等边角碎符。

It was, however, the onset of unemployment which made the worker most conscious of his lowered status in the NEP economy. The continued stagnation of heavy industry, the prices crisis in the consumer industries, the call for rationalization of production, the insistence on khozraschet and on earning profits---all these set up strong pressures for the dismissal of redundant workers. Unemployment resumed its normal role under a market economy as an instrument of labour discipline and pressure on wages. Statistics were few and unreliable. In 1923 the number of unemployed was said to have reached a million; but official returns related to members of trade unions and those registered at the labour exchanges, who were entitled to exiguous relief payments, and took no account of a mass of unskilled workers, mainly peasants, seeking casual work in the towns, especially in the building industry. If NEP had rescued the peasant from disaster, it had reduced industry and the labour market to conditions bordering on chaos. An underground opposition group in the party, which called itself the ``Workers' Group'', and declared that the letters NEP stood for ``new exploitation of the proletariat'', was denounced at the party congress in April. When NEP was freely described as a policy of concessions to the peasant, the question which nobody had asked was at whose expense these concessions were being made. The proletariat, the heroic standard-bearer of the revolution, had suffered dispersal, disintegration and a drastic reduction of numbers under the impact of civil war and industrial chaos. The industrial worker had become the step-child of NEP. 
% 注：去除了 "j\", "llfrole", "/((discipline", "‘ bordering" 的混入标记。

The other crisis, or aspect of the crisis, was financial. The financial consequences of NEP were quite unpremeditated. Once NEP had established the principle of a free market on which goods were bought and sold, these transactions could not be conducted in terms of a constantly declining, and now almost valueless, ruble. The autumn of 1921 saw the introduction of a number of financial reforms. It was decided to draw up the state budget in pre-war rubles, the value of the current ruble being adjusted month by month to this standard. This was in effect a price-index ruble, sometimes referred to as a ``goods'' or ``commodity'' ruble, and was used in the calculation of wage-rates. A State Bank was set up to manage the currency, to re-establish credit, and to lay the foundations of a banking system. At the end of 1921 a party conference advocated the establishment of a currency based on gold; and a few months later the fluctuating ``commodity ruble'' was replaced by a hypothetical ``gold ruble'' as the standard of value. In the autumn of 1922 the State Bank began to issue notes in a new denomination; the chervonets, equivalent to ten gold rubles. But the issue was at first small. For another year the chervonets served as a unit of account, and payments were made in the old paper rubles at the constantly declining current rate. 
% 注：由于 OCR 顺序错乱，"as a unit of aecount, and payments..." 与 "broke in the summer and autumn of 1923. The collapse..." 的顺序被打断。我将前后文完美衔接。

A major economic crisis resulting from all these conditions broke in the summer and autumn of 1923. The collapse of industrial prices in the previous year had impelled the leaders of industry to combine in their own defence. The industrial trusts formed selling syndicates to maintain orderly conditions of marketing and hold up prices. These organizations were remarkably successful in achieving their purpose. By September 1922 the relation between industrial and agricultural prices had been restored to its pre-war balance; and from this time industrial prices rose dramatically at the expense of agricultural prices. Trotsky in his report at the twelfth party congress in April 1923 produced a diagram which showed pictorially how the ``scissors'', representing the blades of agricultural and industrial prices, had opened more and more widely in the past six months. Everyone deplored these violent price fluctuations; but how to prevent them within the framework of NEP was less clear. The party was still deeply committed to the policy of indulgence for the peasant, which was the essence of NEP. Yet the current trend was wholly adverse to the agricultural producer. When in October 1923 the scissors opened to their widest extent, the ratio of industrial prices to agricultural prices was three times as high as in 1913. Meanwhile further monetary problems threatened the economy. To finance the abundant harvest it had been necessary to resume the unlimited printing of ruble notes, thus further depreciating the old paper currency. Attempts were made to substitute the ``gold ruble'' for the ``commodity ruble'' in calculating wage payments; and this was said to reduce actual payments by as much as 40 per cent. This and other grievances of the workers produced a wave of unrest and strikes in the autumn of 1923. 
% 注：去除了 "/leaders", "llconditions", "((were", "^ hj", "/ilfrom", "^lljof", "II923", "mi' t\-vt" 等严重乱码。修正 "priees", "syndieates", "agrieultural", "pietorially", "eeonomy", "ealculating", "reduee" 的 OCR 拼写错误。

The party leaders took alarm at the gathering storm; and the central committee appointed a so-called ``scissors committee'' of 17 members to report on the crisis, with special reference to prices. Trotsky had hitherto been careful not to dissent openly from his colleagues, and perhaps for that reason refused an invitation to serve on the scissors committee. But, while it was deliberating, he lost patience, and on October 8 addressed to the party central committee a letter denouncing ``flagrant radical errors of economic policy''; decisions were being taken without regard to any ``economic plan''. Trotsky condemned ``attempts to command prices in the style of war communism''. The right approach to the peasant was through the proletariat; the rationalization of state industry was the key to the closing of the scissors. The letter was followed a week later by the issue of a ``platform of the 46'', signed by 46 party members, some followers of Trotsky, some of other opposition groups. This spoke of a ``grave economic crisis'', brought on by ``the casual, unconsidered and unsystematic character of the decisions of the central committee''. Both Trotsky's letter and the ``platform'' went on from these criticisms of economic mismanagement to attack the oppressive regime which stifled opinions in the party. 
% 注：删除了 "'''attempts" 中多余的引号。

The ``platform'' had demanded a broad party conference to debate these issues. The central committee responded by opening the columns of Pravda to a controversial discussion---the last of its kind in Soviet history---which lasted for more than a month without participation by any of the leaders, growing more and more confused and heated as it went on. Meanwhile the scissors committee pursued its difficult task. The experience of the past year had convinced almost everyone that prices could not be left to the free play of the market. The committee readily accepted control of wholesale prices. Retail prices presented more difficulty. But it was pointed out that to control wholesale and not retail prices would merely swell the profits of middlemen, identified with the now increasingly unpopular Nepmen. The committee reconciled itself to selective controls of retail prices. But the complexity of the problem, and the timidity of the committee, were such that it did not report till December. 
% 注：去除了 "[of", ".presented", "/control", "4he" 等乱码。

By this time the economic situation had undergone a favourable change. Industrial prices, having reached their peak in October, fell back sharply. The scissors began to close. The harvest, always a major indicator in the primitive Russian economy, was excellent for the second year in succession. Industry, far from being damaged by lower prices, was able to increase its efficiency and enlarge its market. Idle factories and plant were brought back into production. Even pressure on wages was somewhat relaxed. The economic tension of the past six months was eclipsed by a new political tension; this was the moment when the campaign against Trotsky began in earnest. In these circumstances the Politburo adopted a resolution on the report of the scissors committee which was a skilful compromise. The predominance of peasant agriculture was emphasized; nothing must be said to justify Trotsky's insistence on priority for industry. Industry was exhorted to keep prices down, to rationalize itself, and to increase productivity. Control of wholesale prices of articles of mass consumption was to be extended to retail prices: legal maximum prices were to be fixed at once for salt, paraffin and sugar. Concessions were promised on wages, which were to rise ``in proportion to a rise in industry and in the productivity of labour''. Finally there were gestures of support for the financing of heavy industry and for the strengthening of Gosplan. The proposals were endorsed by a party conference in January 1924, a few days before Lenin's death. 
% 注：去除了 "/was", "(were" 等乱码。

The resolution on the report of the scissors committee, for all its caution, gave a certain fillip to industry; by 1924 industry had climbed out of the depths of stagnation and depression in which it floundered when NEP was introduced in 1921. But the revival was one-sided. Light consumer industries directly serving the peasant market prospered. But in NEP conditions nothing stimulated the heavy industries concerned with the production of the means of production, and these lagged behind. According to Gosplan figures, industrial production for the year ended October 1, 1924, while two-and-a-half times as great as in 1920, reached only 40 per cent of the pre-war level, and the metal industries reached only 28.7 per cent. This deficiency began to excite anxiety in the party, notably in opposition circles. The scissors resolution of December 1923 expressed the view that the metal industry should ``be advanced to the front rank and receive from the state support of all kinds''; and this was endorsed by the party conference in January 1924. But nothing was done to implement the pious aspiration. Dzerzhinsky's appointment in February 1924 as president of Vesenkha drew fresh attention to the problem. Three months later Dzerzhinsky reported to the thirteenth party congress that an investment of from 100 to 200 million gold rubles would be required over the next five years to put heavy industry on its feet; and Zinoviev rhetorically exclaimed that it was now ``the turn for metal, the turn for an improvement in the production of means of production, the turn for a revival of heavy industry''. These fine words had no immediate counterpart in action, but they marked a change in the climate of opinion which held out promise for the future. 
% 注：修复了 "October i", "too to 200", "28-7" 的 OCR 错误为 "October 1", "100 to 200", "28.7"。去除了 "i/depression" 的斜杠。

The spring and summer of 1924 were a time of recovery and growing confidence. Agriculture under NEP had emerged from the disasters of the recent past; some indulgence was even shown for the kulak. Industry steadily revived, though the advance was uneven. The currency reform was completed in March 1924, when the gold-based chervonets currency was universally adopted, and the old Soviet ruble notes withdrawn. In May a People's Commissariat of Internal Trade, headed by Kamenev, was set up for the main purpose of operating price controls. The ratio of industrial to agricultural prices had now returned approximately to its 1913 level. Control of industrial prices, both wholesale and retail, seems to have been partially effective, but agricultural prices proved recalcitrant. Foreign trade, managed under the monopoly of foreign trade by a separate commissariat under Krasin, reached sizable dimensions for the first time in the year 1923--1924. Of exports 75 per cent were agricultural products, including grain; the other principal items were timber products and oil. Of imports, nearly 75 per cent were taken by industry, in the form of cotton and other raw materials or semi-manufactured products. These impressive results had been achieved under the regime of NEP, and could not have been achieved without it; they were hailed as a triumphant vindication of NEP. Yet the scissors crisis had been overcome only by measures---especially price controls---which contravened the market principles of NEP; these too had been an essential condition of recovery. Nor was everyone in the party happy about the conspicuous role of the kulak in the villages or of the Nepman in the towns. But the revival in every sector of the economy encouraged the postponement of these baffling problems to a later period. The struggle between the elements of a market economy and of a managed economy went on throughout the nineteen-twenties. 
% 注：清除了 "[in", "Operating" (小写化), "(foreign", "7 oil.", ") in", "V tured", "> r without", "lof" 等等段末的大量碎乱码。